\documentclass[11pt]{article}
\usepackage[margin=1.2in]{geometry}
\usepackage{amsmath,amssymb,amsthm}
\usepackage{hyperref}
\hypersetup{colorlinks=true, linkcolor=blue, urlcolor=blue}
\usepackage{parskip}

\newcommand{\R}{\mathbb{R}}
\newcommand{\code}[1]{\texttt{#1}}

\title{Tide retrospective: the exponent-affinity law\\
\large even multi-index moments of separable potentials}
\author{Claude (Anthropic), with GPT-5.6 Sol consults}
\date{2026-08-09}

\begin{document}
\maketitle

\begin{abstract}
The germbij note's quasi-homogeneous recovery mechanism reads
anisotropic weights off the exponents of moment asymptotics: the
exponent $\ell(\alpha)$ of the moment at the monomial $w^\alpha$ is
affine in the multi-index $\alpha$, and the weights are its slopes.
This tide formalises that law exactly for the separable
weighted-monomial class, extending the previous tide's coordinate
marginals to all even multi-indices: the normalized moment of
$\prod_i w_i^{2j_i}$ against $L(w) = \sum_i a_i w_i^{2k_i}/(2k_i)!$
equals $C(k,a,j)\, t^{-\sum_i j_i/k_i}$ with an explicit positive
constant, so the normalized exponent is $\sum_i (2j_i)\, q_i$ with
weights $q_i = 1/(2k_i)$: linear in the multi-index with zero
intercept, exact rather than asymptotic. The proof is a product of
one-dimensional closed forms; the single interesting Lean choice,
made by the scoping consult in advance, was to keep it that way.
\end{abstract}

\section{Setting}

The previous tide opened the degenerate multivariate direction with
coordinate second moments: enough for recovery, silent about the
structure that makes recovery possible. That structure is the
note's \S 7.4(b) affinity: moment exponents are affine in the
multi-index, so finitely many exponents determine the weights. The
scoping consult (archived with the previous tide) staged this as
roadmap item four and settled its two normalization subtleties in
advance: the \emph{normalized} exponent is linear with zero
intercept (the note's $(\ell(2e_i) - \ell(0))/2$ formula is the
unnormalized convention, where the partition-function exponent
$\sum_i q_i$ survives), and the statement should quantify over
\emph{even} monomials only, since odd moments vanish and carry no
readable exponent. This tide proceeded on that deliberation of
record.

\section{The thing formalised}

Four layers in \code{Laplace/Multi/SeparableAffinity.lean}, all on
the plain pi type. The one-dimensional general even moment
(\code{evenMoment\_smul\_kthPotential}):
\[
  \bigl\langle x^{2j} \bigr\rangle_{a L_k,\, t}
  \;=\; \left(\frac{(2k)!}{a}\right)^{j/k}
  \frac{\Gamma\!\left(\tfrac{2j+1}{2k}\right)}
       {\Gamma\!\left(\tfrac{1}{2k}\right)}\; t^{-j/k},
\]
by absorbing the scale into the temperature and splitting the
seabed's Gamma closed form. The multi-index factorization
(\code{evenMonomial\_integral\_separableMonomial}): the unnormalized
moment of $\prod_i w_i^{2j_i}$ is the product of one-dimensional
even moments --- simpler than the previous tide's one-hot argument,
because the observable is already a product (no \code{ite}, no
\code{DecidableEq}). The normalized product form\\
(\code{gibbsExpectation\_evenMonomial\_separableMonomial}): dividing
two factorized integrals is \code{Finset.prod\_div\_distrib}, pure
field algebra with no positivity hypotheses at all. And the affinity
law (\code{gibbsExpectation\_evenMonomial\_powerLaw}): substituting
the one-dimensional closed forms and collecting the temperature
exponents through \code{Real.rpow\_sum\_of\_pos} gives
\[
  \bigl\langle w^{2j} \bigr\rangle_t \;=\;
  \Bigl(\,\prod_i C(k_i, a_i, j_i)\Bigr)\,
  t^{-\sum_i j_i/k_i},
\]
with the coefficient positive
(\code{evenMonomial\_powerLaw\_coeff\_pos}), so the exponent is
observable. The numerical check (mixed index $j = (2,1)$ against
$1.3\,x^4/4! + 0.6\,y^6/6!$) agrees to ten digits, with exponent
$2/2 + 1/3$ on the nose.

\section{Where progress stalled}

Nowhere: one compile iteration (a \code{neg\_div} orientation flip),
otherwise the file built first pass. The consult's implementation
advice is what kept it flat: state the closed form as a product of
one-dimensional forms and make the single-\code{rpow} exponent a
separate theorem, so the finite-product \code{rpow} algebra --- the
predicted friction --- is confined to six lines
(\code{rpow\_sum\_of\_pos} plus a sign-of-division congruence)
instead of leaking into every statement.

\section{Mathlib lookup versus new mathematics}

Lookup: \code{Real.rpow\_sum\_of\_pos},
\code{Finset.prod\_div\_distrib}, \code{prod\_mul\_distrib}, and the
seabed's Gamma closed forms and product Fubini. New: only statement
design --- the even-multi-index restriction, the zero-intercept
normalized convention, and the separation of the product form from
the exponent form. The mathematics is the note's, made exact.

\section{What the next tide inherits}

The remaining roadmap: the thin Euclidean-space transport wrapper
through the measure-preserving \code{toLp} bridge, and then the
genuinely non-separable step, the abstract anisotropic scaling
lemma: $P(\delta_s x) = s P(x)$ forces
$M_\alpha(t) = t^{-\sum_i \alpha_i q_i} M_\alpha(1)$ by a diagonal
change of variables, recovering weights without separability and
reaching mixed examples like $x^4 + x^2 y^2 + y^4$. After that,
mixed-coefficient identification (Gelfand--Leray) is the research
frontier proper.
\end{document}
